\section{Unified reproducibility and planned artifact release}
\label{sec:unified-reproducibility}

The repository exposes one canonical paper and a tiered verification entry
point.  At the date of this draft, public GitHub and immutable release-asset
publication remain pending.  The commands below specify the reviewed release
contract; they do not assert that an anonymous download already exists.
A clean checkout with Python~3.11.15, exactly
\texttt{uv 0.10.2}, and the locked dependencies is prepared by
\begin{verbatim}
bash scripts/bootstrap.sh
\end{verbatim}
The bootstrap fails closed when that \texttt{uv} version is absent and has no
unhashed package-manager fallback.
The following tiers intentionally make different guarantees.

\begin{center}
\begin{tabularx}{\textwidth}{@{}lXX@{}}
\toprule
Tier & Coverage & Logical scope \\
\midrule
\texttt{quick} & Four current upper checkers, the current lower arithmetic
replay, finite unit/seed checks, and tracked graph certificates. & Fast
integrity check; it does not replay the large finite DRAT proofs. \\
\texttt{full} & \texttt{quick}, the complete long upper interval chain and
the full lower checker suite. & Replays the computer-assisted upper layer and
the lower arithmetic witnesses. \\
\texttt{finite-heavy} & Hash verification, semantic CNF reconstruction, and
six DRAT replays with the pinned checker. & Replays the proof-carrying finite
theorem for the fixed seed. \\
\texttt{paper} & Deterministic materialization and compilation of this single
manuscript. & Checks source integration and typesetting, not the mathematical
certificates by itself. \\
\texttt{sources} & Networked downloads of the three pinned arXiv members and
five upstream graph matrices. & Verifies source identities; it is separate
from the offline theorem replay. \\
\bottomrule
\end{tabularx}
\end{center}

The fast and full tiers are invoked as
\begin{verbatim}
.venv/bin/python reproduce.py quick
.venv/bin/python reproduce.py full
.venv/bin/python reproduce.py sources
\end{verbatim}
Large compressed CNF/DRAT pairs are to be release assets rather than ordinary Git
objects.  After downloading every entry of \path{artifacts/MANIFEST.tsv},
build the source-pinned checker without imposing a cross-platform binary hash:
\begin{verbatim}
bash scripts/build_drat_trim.sh /absolute/new/checker-directory
\end{verbatim}
The checker executable hash is recorded as current-run provenance.  The
proof-carrying tier is
\begin{verbatim}
.venv/bin/python reproduce.py finite-heavy \
  --artifact-dir artifacts/downloads \
  --drat-trim /absolute/path/to/drat-trim
\end{verbatim}
The canonical PDF is generated by \texttt{make paper};
\texttt{make papers} is retained as an alias.  Historical part manuscripts
under \texttt{papers/upper}, \texttt{papers/lower}, and
\texttt{papers/finite} are source components and provenance records, not
separate publication outputs.

The download helper first requires the release asset-name set to equal the
manifest names exactly.  The release manifest then authenticates compressed
artifacts before decompression.
The finite-heavy verifier then reconstructs the structural clauses and
independent-set hitting clauses, checks the exact branch counters, and replays
the proof traces.  This ordering prevents a valid proof of an unrelated or
malformed CNF from being mistaken for evidence about the graph theorem.
